\documentclass[11pt]{article}
\usepackage[paperwidth=220mm,paperheight=320mm,left=16mm,right=16mm,top=18mm,bottom=18mm,headheight=16pt,headsep=10pt,footskip=14pt]{geometry}
\usepackage{fix-cm}
\usepackage{fontspec,amsmath,amssymb,graphicx,array,multirow,tikz}
\usetikzlibrary{decorations.pathreplacing}
\usepackage[unicode,hidelinks]{hyperref}
\usepackage{polyglossia}
\setmainlanguage{latin}
\setotherlanguages{arabic,greek,hebrew,syriac}
\setmainfont{Linux Libertine O}[Ligatures=TeX]
\setsansfont{Linux Biolinum O}[Ligatures=TeX]
\newfontfamily\greekfont{FreeSerif}[Script=Greek]
\newfontfamily\arabicfont{Amiri}[Script=Arabic]
\newfontfamily\hebrewfont{Frank Ruhl Hofshi}[Script=Hebrew]
\newfontfamily\syriacfont{Segoe UI Historic}[Script=Syriac]
\newfontfamily\CapFont{Noto Serif}
\newfontfamily\MoonFont{FreeSerif}
\newcommand{\MoonSym}{\text{\MoonFont ☾}}\newcommand{\SunSym}{\text{\MoonFont ⊙}}
\newfontfamily\EthiopicFont{Ebrima}
\newcommand{\textethiopic}[1]{{\EthiopicFont #1}}
\newfontfamily\NallinoSigns{NallinoSigns.otf}[Path=./fonts/]
\newcommand{\AbjadZero}{{\NallinoSigns\symbol{"E000}}}
\hypersetup{pdftitle={Al-Battani, Opus astronomicum. Nallino, versio Latina et adnotationes (Pars I). S03},pdfauthor={Al-Battani; Carlo Alfonso Nallino (Latin translation and notes)}}
\makeatletter
\def\ps@sourceedition{%
\def\@oddhead{\vbox{\hbox to\textwidth{\fontsize{8}{10}\selectfont C. A. NALLINO — VERSIO ET ADNOTATIONES\hfil\leftmark}\vskip3pt\hrule height.25pt}}%
\let\@evenhead\@oddhead
\def\@oddfoot{\hbox to\textwidth{\fontsize{7}{9}\selectfont\rightmark\hfil\thepage}}%
\let\@evenfoot\@oddfoot}
\makeatother
\pagestyle{sourceedition}
\setlength{\parindent}{0pt}\setlength{\parskip}{0pt}
\newcommand{\BeginSource}[2]{\clearpage\markboth{PAG. #2}{#1}\hypertarget{#1}{}\pdfbookmark[0]{#1 — p. #2}{bk-#1}\fontsize{11.1}{13.8}\selectfont}
\tracinglostchars=2
\newcommand{\Name}[1]{{\addfontfeatures{LetterSpace=4.0}#1}}
\newcommand{\Indent}{\hspace*{1.6em}}
% an Arabic run inside a Latin line: its descenders do not open the line (the printed lines are evenly spaced)
\newcommand{\ArabicRun}[1]{\smash[b]{\textarabic{#1}}}
\newcommand{\qfrac}[2]{{}^{#1}\!/_{{#2}}}
\newcommand{\sa}{\textsuperscript{\textit{a}}}\newcommand{\sm}{\textsuperscript{\textit{m}}}\newcommand{\sd}{\textsuperscript{\textit{d}}}
\newcommand{\RawBlock}[1]{\par\vspace{3pt}\noindent#1\par\vspace{3pt}}
\newcommand{\CenterLine}[1]{\makebox[\linewidth][c]{#1}}
\newlength{\FullSourceWidth}\newlength{\GutterOffset}\newif\ifNumbersLeft
\newcommand{\DSource}[2]{\BeginSource{AB01-PDF#1}{#2}%
 \setlength{\FullSourceWidth}{\textwidth}\setlength{\GutterOffset}{0pt}%
 \ifodd#2\relax\NumbersLefttrue\else\NumbersLeftfalse\fi}
\newsavebox{\DLBox}
\newcommand{\DLine}[4]{%
 \par\noindent\hypertarget{#1}{}%
 \ifx\relax#2\relax\else
   \ifNumbersLeft
    \rlap{\kern-\GutterOffset\llap{\fontsize{7}{8}\selectfont #2\hspace{3mm}}}%
   \else
    \rlap{\kern\dimexpr\FullSourceWidth-\GutterOffset\relax\hspace{3mm}{\fontsize{7}{8}\selectfont #2}}%
   \fi
 \fi
 \ifx\relax#3\relax\else
   \ifNumbersLeft
    \rlap{\kern\dimexpr\FullSourceWidth-\GutterOffset\relax\hspace{3mm}{\fontsize{7}{8}\selectfont #3}}%
   \else
    \rlap{\kern-\GutterOffset\llap{\fontsize{7}{8}\selectfont #3\hspace{3mm}}}%
   \fi
 \fi
 \sbox{\DLBox}{\strut #4}%
 \ifdim\wd\DLBox>\linewidth\typeout{DIPLOMATICWIDTH|#1|\the\wd\DLBox|\the\linewidth}\fi
 \usebox{\DLBox}\strut\par}
\begin{document}
\thispagestyle{empty}
\begin{center}\Large AL-BATTĀNĪ\par\vspace{4mm}\LARGE OPUS ASTRONOMICUM\par
\vspace{5mm}\large Caroli Alphonsi Nallino versio Latina et adnotationes (pars I)\par\vspace{8mm}
\Large S03 --- editio diplomatica\par\vspace{4mm}\large Paginae impressae 110--110\end{center}
\vspace{12mm}\noindent Singulae lineae paginarum impressarum singulis lineis huius editionis respondent, cum
divisionibus vocabulorum, signis et generibus litterarum (rectis, inclinatis, distantibus). Numeri linearum et
paginae textus Arabici in marginibus stant ut impressi sunt. Textus ex imaginibus paginarum transcriptus est;
litterae Graecae et Arabicae, numeri et signa in imaginibus amplificatis lecta sunt.
Haec transcriptio a nullo homine recognita est.\par
\DSource{0199}{110}
\DLine{AB01-PDF0199-body-L001}{}{}{latitudinis Lunae numera; ad gradum horizontis ad quem ita pervenies, obscuritas tendet.}
\DLine{AB01-PDF0199-body-L002}{}{}{Scilicet, exstante Sole ab occidente, numerare incipies ab azimuth gradus occidentis; exstante}
\DLine{AB01-PDF0199-body-L003}{}{}{vero Sole in parte orientali, ab azimuth gradus ascendentis inchoabis.}
\hypertarget{AB01-PDF0199-P01}{}
\DLine{AB01-PDF0199-body-L004}{}{}{\Indent Si per tabulas, modo veritati propinquo, eclipses Solis cognoscere vis, cum Solem deficere}
\DLine{AB01-PDF0199-body-L005}{5}{}{scias aut saltem deficere posse, quantum in tua urbe horae aequinoctiales verae coniunctionis}
\DLine{AB01-PDF0199-body-L006}{}{}{a tempore meridiano distent hoc modo deprehende. Horas coniunctionis in horas tui loci con-}
\DLine{AB01-PDF0199-body-L007}{}{}{verte; eas de 24 horis deme, si coniunctio est antemeridiana, aut eas serva, si coniunctio post}
\DLine{AB01-PDF0199-body-L008}{}{}{meridiem evenit usque ad occasum Solis. Habebis distantiam horarum coniunctionis a linea}
\DLine{AB01-PDF0199-body-L009}{}{}{meridiana. Per eam in tabulis parallaxium Lunae in climatibus (1) cognosce parallaxim lunarem}
\DLine{AB01-PDF0199-body-L010}{10}{}{longitudinis, ut postea scias parallaxim correctam ad gradum in quo Luna reperitur per quar-}
\DLine{AB01-PDF0199-body-L011}{}{}{tam (2) columnam tabulae correctionis (3). Parallaxim correctam divide per motum horarium}
\DLine{AB01-PDF0199-body-L012}{}{}{Lunae; exibunt horae primae parallaxis, quas ab horis verae coniunctionis deme vel iis adde,}
\DLine{AB01-PDF0199-body-L013}{}{}{prout Luna minus vel magis quam 90° ab ascendente distat. — Earum distantiam a meridie}
\DLine{AB01-PDF0199-body-L014}{}{}{iterum sume; per eam deprehende parallaxim Lunae in longitudine, correctam quarta (4) co-}
\DLine{AB01-PDF0199-body-L015}{15}{p. 167.}{lumna tabulae correctionis; haec erit secunda parallaxis, quam per praecessionem Lunae [re-}
\DLine{AB01-PDF0199-body-L016}{}{}{spectu Solis] divide, ut horas secundae parallaxis habeas. Eas ab horis verae coniunctionis}
\DLine{AB01-PDF0199-body-L017}{}{}{deme, aut iis adde, prout distantia coniunctionis scilicet gradus Lunae ab ascendente minor}
\DLine{AB01-PDF0199-body-L018}{}{}{aut maior quam 90° est. — Earum distantiam a meridie tertia vice sume; per eam cognosce}
\DLine{AB01-PDF0199-body-L019}{}{}{parallaxim correctam gradus Lunae in longitudine; quae si ut secunda parallaxis est, horae}
\DLine{AB01-PDF0199-body-L020}{20}{}{verae coniunctionis, post additas vel demptas horas secundae parallaxis, sunt horae medii ecli-}
\DLine{AB01-PDF0199-body-L021}{}{}{psis. — At si tertia parallaxis secundam superat, earum differentiam observa. Postea sextantem}
\DLine{AB01-PDF0199-body-L022}{}{}{horae adde horis distantiae coniunctionis a meridie, quas per secundam parallaxim inveneras;}
\DLine{AB01-PDF0199-body-L023}{}{}{hac summa cognita, inveni, ut soles, parallaxim longitudinis correctam gradus Lunae. Quan-}
\DLine{AB01-PDF0199-body-L024}{}{}{titatem qua haec superat secundam parallaxim, multiplica in sex, et productum de praecessione}
\DLine{AB01-PDF0199-body-L025}{25}{}{horaria Lunae [respectu Solis] (5) deme; remanet motus aequatus, per quem divide \textit{differen-}}
\DLine{AB01-PDF0199-body-L026}{}{}{\textit{tiam inter tertiam et secundam parallaxim, quam te servare iussi} (6). Exibit pars horae;}
\DLine{AB01-PDF0199-body-L027}{}{}{eam adde horis secundae parallaxis, et habebis horas praecisas (7). — Denique si tertia paralla-}
\DLine{AB01-PDF0199-body-L028}{}{}{xis est secunda minor, earum differentiam in 6 multiplicatam adde praecessioni Lunae; habebis}
\DLine{AB01-PDF0199-body-L029}{}{}{motum aequatum, per quem divide differentiam duarum parallaxium. Exibit pars horae, quam}
\DLine{AB01-PDF0199-body-L030}{30}{}{ab horis secundae parallaxis deme, ut horas praecisas habeas.}
\hypertarget{AB01-PDF0199-P02}{}
\DLine{AB01-PDF0199-body-L031}{}{}{\Indent Horas praecisas parallaxis separatim multiplica in motum inaequalem Lunae et motum}
\DLine{AB01-PDF0199-body-L032}{}{}{inaequalem Solis in ea hora; ambo producta serva. Si tunc Luna ab ascendente minus quam}
\DLine{AB01-PDF0199-body-L033}{}{}{90° distat, horas praecisas secundae parallaxis deme de horis verae coniunctionis, motum Solis}
\DLine{AB01-PDF0199-body-L034}{}{}{eo [temporis] spatio de loco coniunctionis, denique motum Lunae de loco coniunctionis, de}
\DLine{AB01-PDF0199-body-L035}{35}{}{anomalia Lunae et de motu latitudinis; at de motu latitudinis eum cum motu nodi [borealis]}
\DLine{AB01-PDF0199-body-L036}{}{}{in horis secundae parallaxis detrahe. Si contra Luna ab ascendente magis quam 90° distat,}
\DLine{AB01-PDF0199-body-L037}{}{}{additionem pro subductione adhibe. — Habebis horas medii eclipsis, locum Solis, locum Lu-}
\DLine{AB01-PDF0199-body-L038}{}{}{nae, anomaliam Lunae et motum latitudinis, omnia correcta ad medium eclipsis.}
\par\vspace{6pt}\nointerlineskip\hbox to\textwidth{\hss\vrule width 115.43pt height.25pt depth0pt\hss}\prevdepth=-1000pt \vspace{5pt}
\noindent\begin{minipage}[t]{.48\textwidth}\vspace{0pt}\fontsize{9}{11.5}\selectfont\setlength{\GutterOffset}{0pt}\setlength{\FullSourceWidth}{\linewidth}
\hypertarget{AB01-PDF0199-N01}{}
\DLine{AB01-PDF0199-notes_left-L001}{}{}{\Indent (1) Fol. 198,v. et seqq.}
\hypertarget{AB01-PDF0199-N02}{}
\DLine{AB01-PDF0199-notes_left-L002}{}{}{\Indent (2) Ita cod. et Plato, sed addendum « et quin-}
\DLine{AB01-PDF0199-notes_left-L003}{}{}{« tam » ut ex regulis cap. XXXIX, pag. 83 patet.}
\hypertarget{AB01-PDF0199-N03}{}
\DLine{AB01-PDF0199-notes_left-L004}{}{}{\Indent (3) Fol. 195,r. — Quae usque ad pag. 111, l. 6,}
\DLine{AB01-PDF0199-notes_left-L005}{}{}{sequuntur sunt iam uberius in pag. 104, l. 21—p. 106,}
\DLine{AB01-PDF0199-notes_left-L006}{}{}{l. 22 exposita. Sed auctor nunc vult tabularum usum}
\DLine{AB01-PDF0199-notes_left-L007}{}{}{explicare.}
\end{minipage}
\hfill\vrule width.3pt\hfill\begin{minipage}[t]{.48\textwidth}\vspace{0pt}\fontsize{9}{11.5}\selectfont\setlength{\GutterOffset}{0pt}\setlength{\FullSourceWidth}{\linewidth}
\hypertarget{AB01-PDF0199-N04}{}
\DLine{AB01-PDF0199-notes_right-L001}{}{}{\Indent (4) Adde, ut supra dixi, « et quinta ».}
\hypertarget{AB01-PDF0199-N05}{}
\DLine{AB01-PDF0199-notes_right-L002}{}{}{\Indent (5) I. e. de motu synodico (vel relativo) horario}
\DLine{AB01-PDF0199-notes_right-L003}{}{}{Lunae.}
\hypertarget{AB01-PDF0199-N06}{}
\DLine{AB01-PDF0199-notes_right-L004}{}{}{\Indent (6) Vide quae de locis parallelis pag. 105, l. 24}
\DLine{AB01-PDF0199-notes_right-L005}{}{}{et 34, disseritur in adnot. \textbf{\textit{a}}.}
\hypertarget{AB01-PDF0199-N07}{}
\DLine{AB01-PDF0199-notes_right-L006}{}{}{\Indent (7) Vel aequatas, ut in loco parallelo pag. 105,}
\DLine{AB01-PDF0199-notes_right-L007}{}{}{l. 37 sqq., legitur.}
\end{minipage}
\end{document}
